\SourceLeaf{290}%
\[
\begin{aligned}
\mathrm{d}\mathord{.}\frac{\operatorname{tang}\psi}{m}&=-\frac{l\mathrm{d}x}{m}\\
\mathrm{d}u&=\operatorname{tang}\psi\mathord{.}\mathrm{d}x
\end{aligned}
\]
angewandt, so ergibt sich
\[
\begin{aligned}
\operatorname{tang}\psi&=\mathfrak{A}(1+\mu x^3+\mu'x^4+\text{u.\,s.\,w.})-\tfrac{1}{3}\lambda x^3-\tfrac{1}{4}\lambda'x^4-\text{u.\,s.\,w.}\\
u&=\mathfrak{B}+\mathfrak{A}(x+\tfrac{1}{4}\mu x^4+\tfrac{1}{5}\mu'x^5+\text{u.\,s.\,w.})-\tfrac{1}{12}\lambda x^4-\tfrac{1}{20}\lambda'x^5-\text{u.\,s.\,w.}
\end{aligned}
\]
\indent Die durch die Integration eingeführten Constanten, \(\mathfrak{A},\mathfrak{B}\), lassen sich durch die Bedingung bestimmen, dass \(u=0\) werden muss für die beiden Werthe von \(x\), welche den Punkten \(F,G\) entsprechen. Es seien diese Werthe \(x=-\tfrac{1}{2}(h-\delta)\) und \(x=+\tfrac{1}{2}(h+\delta)\), wo \(\delta\) den Werth von \(2x\) in dem mitten zwischen \(F\) und \(G\) liegenden Punkte ausdrückt, und allgemein zu reden eine Grösse von derselben Ordnung wie \(h\) ist, oder von einer höhern, wenn \(H\) dieser Mitte sehr nahe liegt. Man leitet hieraus leicht folgenden auf die Ordnung \(h^3\) (einschl.) genauen Ausdruck für \(\mathfrak{A}\) ab
\[
\mathfrak{A}=\frac{\lambda((h+\delta)^4-(h-\delta)^4)}{192h}=\tfrac{1}{24}\lambda\delta(hh+\delta\delta)
\]
\indent Substituirt man diesen in der Reihe für \(\operatorname{tang}\psi\). und legt dann der Veränderlichen \(x\) die bestimmten Werthe \(-\tfrac{1}{2}(h-\delta),+\tfrac{1}{2}(h+\delta)\) bei, so ergibt sich, gleichfalls auf die dritte Ordnung genau,
\[
\begin{aligned}
\operatorname{tang}\psi^0&=\tfrac{1}{24}\lambda h(hh-2h\delta+3\delta\delta)\\
\operatorname{tang}\psi'&=-\tfrac{1}{24}\lambda h(hh-2h\delta+3\delta\delta)
\end{aligned}
\]
\indent In dem speciellen Fall der in der Folge zu entwickelnden Untersuchung kommt übrigens zu der oben bezeichneten Bedingung noch der Umstand hinzu, dass der Normalparallelkreis mitten inne liegt zwischen den beiden Parallelkreisen, auf welchen sich die Punkte \(F,G\) befinden, und in Folge dieses Umstandes werden schon die abgekürzten Ausdrücke
\[
\begin{aligned}
\operatorname{tang}\psi^0&=\tfrac{1}{24}\lambda h^3\\
\operatorname{tang}\psi'&=-\tfrac{1}{24}\lambda h^3
\end{aligned}
\]
auf die dritte Ordnung genau sein, wie sich leicht auf folgende Art zeigen lässt. Bezeichnet man die Breite von \(F\) mit \(Q+q\), die von \(G\) mit \(Q-q\), so geben die sphaerischen Dreiecke \(F,H\), Pol und \(G,H\), Pol die Gleichungen
